\section{Introduction}
\label{sec:introduction}

\subsection{Facilitation and persistence}

We consider spin systems in which one state facilitates updates at neighbouring
sites. We label this state~$0$ and the other state~$1$, following the convention
for kinetically constrained spin models (KCSM)
\cite{CancriniEtAl2008,HartarskyToninelli2025}. The interpretation depends on
the model: a~$0$ may infect a neighbour, permit a constrained refresh, or cause
a neighbouring spin to flip. In the facilitating part of the dynamics, a
region containing only~$1$'s has no internal source of constrained updates,
whereas an update enabled by a~$0$ may create further~$0$'s. The two states
therefore play different roles. This does not imply monotonicity of the
dynamics. A refresh or flip may create or remove a~$0$, and ordered
configurations need not remain ordered.

We ask whether this asymmetry yields a comparison of monomial moments and
whether independent~$1\to0$ updates erase their dependence on the initial law.
The main examples do not preserve the product order on configurations, and an
independent perturbation need not preserve a reversible measure. We therefore
seek an argument based on the graphical construction which neither orders
trajectories nor requires the limiting measure to be known.

The effect of a facilitating update need not persist until it is transmitted to
another site. If no information leaves a site between two local updates, a
second birth is redundant, a later refresh overwrites the earlier refresh, and
two flips cancel. For a monomial
$\chi_A(\eta)=\prod_{i\in A}\eta(i)$, the signed Feynman--Kac dual follows the
dependence of~$P_t\chi_A$ backward through the graphical construction. Its signs
record cancellations among local updates. Taking absolute values after each
dual update loses these cancellations. The local updates should instead be
averaged until the next interaction which transmits dependence between sites.

\subsection{Patches and main results}

Starting the dual from~$A$, we retain the times, sources, and target sets of the
interactions which occur at active sources and have nonempty target. Their
records form the successful-interaction skeleton. The skeleton partitions each
site line into intervals, called patches, between successive incoming or
outgoing successful interactions. Conditional on the skeleton, the remaining
marked interactions in distinct patches are independent. Averaging the signed
Feynman--Kac factor on each patch gives the exact semigroup representation in
Theorem~\ref{thm:patch-representation}. In this representation, consecutive
local updates are averaged before their contributions are compared.

The averaged contribution of a completed patch~$P$ is denoted by~$C(P)$. Patch
positivity is the condition~$C(P)\geq0$ for every patch shape.
Theorem~\ref{thm:patch-positivity-order} characterizes this condition by
inequalities for the multilinear coefficients of the rates~$0\to1$ and
$1\to0$. The same coefficients determine the smallest terminal one-density
profile~$\mb p^\star$ for which every end-patch contribution is nonnegative.
This is a threshold for the sign argument, not a phase transition or an
ergodicity threshold.

Products centered at~$\mb p^\star$ define an order~$\preceq_*$ on probability
measures. Theorem~\ref{thm:patch-positivity-order} shows that the semigroup of a
patch-positive system preserves this order, even though it need not preserve
the pointwise order of configurations. The proof is algebraic: in finite volume
the generator has nonnegative off-diagonal entries in the centered-monomial
basis.
Corollary~\ref{cor:pure-death-comparison} compares the system with the dynamics
obtained by removing an environment-independent~$1\to0$ component.

Suppose now that the interaction graph has polynomial growth of exponent~$D$
and that the dynamics contains independent~$1\to0$ updates at rate
$\varepsilon>0$. Theorem~\ref{thm:common-invariant-limit} shows that every
initial law in the class~$\mathcal M_-$ converges to the same invariant measure,
with local convergence rate

\begin{equation*}
O\bb{(1+t)^D e^{-\varepsilon t/2}}.
\end{equation*}

The class~$\mathcal M_-$ is an affine extension of the measures with
nonnegative centered moments. If~$\mb p^\star\leq\mbs{\tfrac12}$, it contains
every probability measure, so the invariant measure is unique and convergence
is uniform over initial configurations. The proof cuts the dual construction
at an intermediate time. Finite propagation controls the event that the
earlier skeleton leaves a linearly growing directed ball; the independent
$1\to0$ component controls both successful interactions after the cut and the
remaining dependence of the end patches on the terminal configuration.

\subsection{Relation to previous work}

Set-valued graphical constructions and additive dual processes go back to
Harris~\cite{Harris1974,Harris1978}; see also
Griffeath~\cite{Griffeath1979}, Holley--Stroock~\cite{HolleyStroock1979}, and
Gray~\cite{Gray1986}. We use the standard monomial Feynman--Kac duality under
the usual finite-speed assumptions for spin systems; Jansen and
Kurt~\cite{JansenKurt2014} give a general account of Markov-process duality.
Conditioning on the successful skeleton gives the one-site factorization over
patches which is not present in the standard duality. Information percolation
also retains updates which carry information~\cite{LubetzkySly2016}, but its
clusters and oblivious updates are used for a different decomposition. The
patch factorization retains cancellations among local updates. It is also
different from the time-scaling observation
in~\cite{GluchowskiMenz2025}, where each local update transition matrix is mixed
with the identity. The construction was motivated by the wall-type one-sided
spin systems studied in~\cite{GluchowskiMenz2026}; many of their rate functions
satisfy the coefficient criterion below.

Classical correlation inequalities for spin systems are tied to attractiveness
and association~\cite{Harris1977,Liggett2006}; duality can also yield moment
inequalities in particular models~\cite{GiardinaRedigVafayi2010}. The order
considered here uses products centered at~$\mb p^\star$ and applies to
non-attractive systems. To our knowledge, preservation of this centered-moment
order has not previously been obtained from a local coefficient condition of
the present form. The convergence theorem is also distinct from classical
positive-rates results such as Gray~\cite{Gray1982}: it assumes patch positivity
and a specified independent component and proves convergence from the
class~$\mathcal M_-$.

The examples come from two lines of work. FA-1f originates
in~\cite{FredricksonAndersen1984} and is placed in the KCSM framework
in~\cite{CancriniEtAl2008,HartarskyToninelli2025}. Its out-of-equilibrium
theory, and that of related KCSMs such as East, is developed
in~\cite{CancriniEtAl2010,BlondelEtAl2013,MountfordValle2019,Mareche2019,
HartarskyToninelli2024}. BABP was introduced by Neuhauser and
Sudbury~\cite{NeuhauserSudbury1993}; its qualitative convergence and
quasi-duality theory are due in part to Sudbury and
Lloyd--Sudbury~\cite{Sudbury1997,LloydSudbury1997}. The product-law estimates
in~\cite{MartinelliShapiraToninelli2025} provide the endpoints for the
centered-moment corollary proved here. These examples show how the same
coefficient criterion applies to different facilitation rules.

\subsection{Organization}

Section~\ref{sec:setup} fixes the notation, and
Section~\ref{sec:main-results} states the three principal theorems.
Section~\ref{sec:signed-dual} constructs the successful skeleton and proves the
patch representation, Theorem~\ref{thm:patch-representation}.
Section~\ref{sec:patch-positivity} proves the coefficient criterion and the
centered-moment comparisons in Theorem~\ref{thm:patch-positivity-order}, as
well as Corollary~\ref{cor:pure-death-comparison}.
Section~\ref{sec:convergence} proves the common invariant-limit theorem.
The FA-1f and BABP calculations are collected in
Section~\ref{sec:applications}. Generator calculations, one-site formulas, and
propagation estimates are deferred to the appendices.
